% --- CORRECCIÓN IMPORTANTE EN LA PRIMERA LÍNEA ---
% Agregamos "xcolor=table" para que funcione \rowcolor sin errores
\documentclass[aspectratio=169, 10pt, xcolor=table]{beamer}

% --- Configuración del Tema (estilo profesional como el ejemplo) ---
\usetheme{Madrid}
\usecolortheme{dolphin}

% --- Paquetes necesarios ---
\usepackage[utf8]{inputenc}
\usepackage[spanish]{babel}
\usepackage{amsmath, amssymb}
\usepackage{booktabs}
\usepackage{graphicx}
\usepackage{listings}
\usepackage{tikz}
\usetikzlibrary{arrows.meta, positioning, shapes.geometric}

% --- Ruta de Imágenes (crea la carpeta "imagenes" y coloca ahí tus archivos) ---
\graphicspath{{./imagenes/}}

% --- Estilo para código Python (configurado para evitar problemas con UTF-8) ---
\definecolor{codegreen}{rgb}{0,0.6,0}
\definecolor{codegray}{rgb}{0.5,0.5,0.5}
\definecolor{codepurple}{rgb}{0.58,0,0.82}
\definecolor{backcolour}{rgb}{0.95,0.95,0.92}
\lstset{
    language=Python,
    backgroundcolor=\color{backcolour},
    commentstyle=\color{codegreen},
    keywordstyle=\color{blue},
    numberstyle=\tiny\color{codegray},
    stringstyle=\color{codepurple},
    basicstyle=\ttfamily\scriptsize,
    breaklines=true,
    frame=single,
    showstringspaces=false,
    literate={á}{{\'a}}1 {é}{{\'e}}1 {í}{{\'i}}1 {ó}{{\'o}}1 {ú}{{\'u}}1 {ñ}{{\~n}}1 {ü}{{\"u}}1
}

% --- Pie de página personalizado ---
\makeatletter

\setbeamertemplate{footline}{
  \leavevmode%
  \hbox{%
  \begin{beamercolorbox}[wd=.333333\paperwidth,ht=2.25ex,dp=1ex,center]{author in head/foot}%
    \usebeamerfont{author in head/foot}\insertshortauthor
  \end{beamercolorbox}%
  \begin{beamercolorbox}[wd=.333333\paperwidth,ht=2.25ex,dp=1ex,center]{title in head/foot}%
    \usebeamerfont{title in head/foot}Sesi\'on 12
  \end{beamercolorbox}%
  \begin{beamercolorbox}[wd=.333333\paperwidth,ht=2.25ex,dp=1ex,right]{date in head/foot}%
    \usebeamerfont{date in head/foot}NLP \hfill \insertframenumber/\inserttotalframenumber \hspace*{0.3cm}
  \end{beamercolorbox}}%
  \vskip0pt%
}

\makeatother

% --- Datos de la portada ---
\title[SESI\'ON 12 NLP]{\textbf{Sesi\'on 12}}
\subtitle{Retrieval-Augmented Generation (RAG) y Bases de Datos Vectoriales}
\author[Prof. C. S\'anchez]{Carlos C\'esar S\'anchez Coronel}
\institute{\textbf{NLP}}
\date{}

% Logo opcional (descomentar si tienes la imagen)
% \titlegraphic{\includegraphics[height=1.5cm]{logo.png}}

% --- Eliminar la barra de navegación (opcional) ---
\setbeamertemplate{navigation symbols}{}

% --- Índice al inicio de cada sección ---
\AtBeginSection[]{
  \begin{frame}
    \frametitle{Contenido}
    \tableofcontents[currentsection]
  \end{frame}
}

% ============================================================================
% INICIO DEL DOCUMENTO
% ============================================================================
\begin{document}

% --- Portada ---
\begin{frame}
    \titlepage
\end{frame}

% --- Índice general ---
\begin{frame}{Contenido}
    \tableofcontents
\end{frame}

% ============================================================================
% SECCIÓN 1: LIMITACIONES DE LOS LLMS
% ============================================================================
\section{Limitaciones de los LLMs}

\begin{frame}{Problemas de los modelos generativos puros}
    \begin{itemize}
        \item \textbf{Conocimiento desactualizado}: entrenados con un corte temporal (ej. 2023). No saben eventos recientes.
        \item \textbf{Alucinaciones}: generan respuestas falsas pero coherentes cuando no tienen informaci\'on.
        \item \textbf{Datos privados}: no conocen documentos internos de una empresa.
        \item \textbf{Soluci\'on}: Retrieval-Augmented Generation (RAG).
    \end{itemize}
\end{frame}

% ============================================================================
% SECCIÓN 2: ¿QUÉ ES RAG?
% ============================================================================
\section{¿Qu\'e es RAG?}

\begin{frame}{Definici\'on de RAG}
    \begin{block}{RAG (Retrieval-Augmented Generation)}
        Framework que optimiza la salida de un LLM consultando una base de conocimiento externa antes de generar la respuesta.
    \end{block}
    \begin{columns}[t]
        \column{0.33\textwidth}
        \textbf{Retrieval}
        \begin{itemize}
            \item Busca documentos relevantes en una base de datos.
        \end{itemize}
        \column{0.33\textwidth}
        \textbf{Augmentation}
        \begin{itemize}
            \item Combina consulta y documentos en un prompt enriquecido.
        \end{itemize}
        \column{0.33\textwidth}
        \textbf{Generation}
        \begin{itemize}
            \item LLM genera respuesta basada en el prompt.
        \end{itemize}
    \end{columns}
\end{frame}

\begin{frame}{Ventajas de RAG}
    \begin{itemize}
        \item \textbf{Conocimiento actualizado}: se actualiza la base, no el modelo.
        \item \textbf{Reducci\'on de alucinaciones}: el modelo se fuerza a basarse en el contexto.
        \item \textbf{Datos privados}: no se exponen en el entrenamiento.
        \item \textbf{Transparencia}: puede citar fuentes.
    \end{itemize}
\end{frame}

% ============================================================================
% SECCIÓN 3: EMBEDDINGS Y BÚSQUEDA SEMÁNTICA
% ============================================================================
\section{Embeddings y B\'usqueda Sem\'antica}

\begin{frame}{Embeddings de texto}
    \begin{itemize}
        \item Representaciones vectoriales densas de fragmentos de texto.
        \item Textos con significado similar tienen vectores cercanos en el espacio.
        \item Ejemplo: $ \text{Vector}(\text{``Rey''}) - \text{Vector}(\text{``Hombre''}) + \text{Vector}(\text{``Mujer''}) \approx \text{Vector}(\text{``Reina''}) $
        \item Modelos: Sentence-BERT, E5, BGE, OpenAI embeddings.
    \end{itemize}
\end{frame}

\begin{frame}{Similitud del coseno}
    \begin{block}{F\'ormula}
        \[
        \text{sim}(\mathbf{A}, \mathbf{B}) = \cos(\theta) = \frac{\mathbf{A} \cdot \mathbf{B}}{\|\mathbf{A}\| \|\mathbf{B}\|}
        \]
        \begin{itemize}
            \item Mide el \'angulo entre dos vectores.
            \item Valores cercanos a 1 indican alta similitud sem\'antica.
            \item Es la m\'etrica est\'andar en bases vectoriales.
        \end{itemize}
    \end{block}
\end{frame}

% ============================================================================
% SECCIÓN 4: BASES DE DATOS VECTORIALES
% ============================================================================
\section{Bases de Datos Vectoriales}

\begin{frame}{Tecnolog\'ias de almacenamiento vectorial}
    \begin{table}
        \centering
        \begin{tabular}{l|c|c}
            \toprule
            \textbf{Herramienta} & \textbf{Tipo} & \textbf{Uso Principal} \\
            \midrule
            FAISS & Librer\'ia (Meta) & B\'usqueda local, alta velocidad en RAM. \\
            Pinecone & Servicio gestionado & Escalabilidad en la nube, nivel empresarial. \\
            Weaviate & Open Source & Base de datos con b\'usqueda h\'ibrida. \\
            ChromaDB & Open Source & F\'acil integraci\'on con Python, prototipos. \\
            \bottomrule
        \end{tabular}
    \end{table}
\end{frame}

% ============================================================================
% SECCIÓN 5: EVALUACIÓN DE SISTEMAS RAG (RAG TRIAD)
% ============================================================================
\section{Evaluaci\'on de Sistemas RAG}

\begin{frame}{El RAG Triad}
    \begin{itemize}
        \item \textbf{Faithfulness (Fidelidad)}: ¿La respuesta se deriva estrictamente del contexto recuperado? (no alucina)
        \item \textbf{Relevance (Relevancia del contexto)}: ¿Los documentos recuperados son \'utiles para responder la pregunta?
        \item \textbf{Answer Relevance}: ¿La respuesta responde la pregunta del usuario? (puede medirse con similitud sem\'antica)
    \end{itemize}
\end{frame}

\begin{frame}{M\'etricas offline para retrieval}
    \begin{itemize}
        \item \textbf{Recall@k}: proporci\'on de documentos relevantes recuperados en top-$k$.
        \item \textbf{MRR (Mean Reciprocal Rank)}: posici\'on del primer documento relevante.
        \item \textbf{NDCG}: relevancia con descuento por posici\'on.
        \item Se necesita un dataset de evaluaci\'on con pares pregunta-documento relevante.
    \end{itemize}
\end{frame}

% ============================================================================
% SECCIÓN 6: IMPLEMENTACIÓN CONCEPTUAL EN PYTHON
% ============================================================================
\section{Implementaci\'on Conceptual}

\begin{frame}[fragile]{Ejemplo con LangChain}
    \begin{lstlisting}
from langchain_community.vectorstores import FAISS
from langchain_openai import OpenAIEmbeddings, ChatGPT
from langchain.chains import RetrievalQA

# 1. Preparar datos y embeddings
embeddings = OpenAIEmbeddings()
texts = ["La Sesion 12 trata sobre RAG", "Carlos es el profesor"]
vector_db = FAISS.from_texts(texts, embeddings)

# 2. Configurar el LLM
llm = ChatGPT(model_name="gpt-4", temperature=0)

# 3. Crear la cadena RAG
rag_chain = RetrievalQA.from_chain_type(
    llm=llm,
    chain_type="stuff",
    retriever=vector_db.as_retriever()
)

# 4. Consulta
query = "De que trata la Sesion 12?"
response = rag_chain.invoke(query)
print(response["result"])
    \end{lstlisting}
\end{frame}

% ============================================================================
% SECCIÓN 7: RAG vs FINE-TUNING
% ============================================================================
\section{RAG vs Fine-Tuning}

\begin{frame}{Comparaci\'on RAG vs Fine-Tuning}
    \begin{table}
        \centering
        \begin{tabular}{l|c|c}
            \toprule
            \textbf{Caracter\'istica} & \textbf{Fine-Tuning} & \textbf{RAG} \\
            \midrule
            Conocimiento & Est\'atico (en pesos) & Din\'amico (base de datos) \\
            Costo & Alto (requiere GPUs) & Bajo (solo inferencia) \\
            Alucinaciones & Posibles & Reducidas \\
            Transparencia & Baja & Alta (cita fuentes) \\
            Actualizaci\'on & Reentrenar & Indexar nuevos docs \\
            \bottomrule
        \end{tabular}
    \end{table}
\end{frame}

% ============================================================================
% SECCIÓN 8: PREPARACIÓN PARA ENTREVISTAS
% ============================================================================
\section{Preparaci\'on para Entrevistas}

\begin{frame}{Preguntas t\'ipicas}
    \begin{itemize}
        \item ¿Qu\'e es RAG y cu\'ando lo usar\'ias?
        \item Diferencias entre RAG y Fine-Tuning.
        \item ¿Qu\'e son los embeddings y la similitud del coseno?
        \item ¿Qu\'e bases de datos vectoriales conoces?
        \item ¿C\'omo evaluar\'ias un sistema RAG? (RAG Triad)
    \end{itemize}
\end{frame}

\begin{frame}{Preguntas avanzadas}
    \begin{itemize}
        \item ¿C\'omo elegir\'ias el tama\~no de chunk para los documentos?
        \item ¿Qu\'e estrategias de re-ranking propondr\'ias?
        \item ¿C\'omo manejar\'ias consultas que requieren informaci\'on de m\'ultiples documentos?
        \item ¿Qu\'e consideraciones de privacidad hay al usar RAG en salud?
        \item ¿C\'omo implementar\'ias un sistema RAG multimodal?
    \end{itemize}
\end{frame}

% ============================================================================
% SECCIÓN 9: CASO INTEGRADO - SISTEMA DE INTELIGENCIA DE NEGOCIOS
% ============================================================================
\section{Caso Integrado: Inteligencia de Negocios en Importaci\'on/Exportaci\'on}

\begin{frame}{Problema}
    \begin{itemize}
        \item Empresa con miles de documentos: regulaciones aduaneras, aranceles, tratados, contratos.
        \item Informaci\'on dispersa en PDFs, bases de datos, sitios web que cambian frecuentemente.
        \item Analistas necesitan respuestas r\'apidas a preguntas como ``¿Cu\'al es el arancel para importar textiles desde Vietnam?''
    \end{itemize}
\end{frame}

\begin{frame}{Arquitectura propuesta}
    \begin{enumerate}
        \item \textbf{Procesamiento}: extraer texto de PDFs, dividir en chunks sem\'anticos (por art\'iculos). Convertir tablas a texto.
        \item \textbf{Embeddings}: usar BGE-M3 o E5 (multiling\"ue, alta calidad).
        \item \textbf{Base vectorial}: Weaviate (open source, b\'usqueda h\'ibrida, filtrado por metadatos).
        \item \textbf{LLM}: modelo local (Llama 3, Mistral) por privacidad y costo.
    \end{enumerate}
\end{frame}

\begin{frame}{Actualizaci\'on y evaluaci\'on}
    \begin{itemize}
        \item \textbf{Indexaci\'on incremental}: al llegar nuevos documentos, se generan embeddings y se insertan.
        \item \textbf{Evaluaci\'on offline}: dataset de preguntas con documentos relevantes anotados; m\'etricas Recall@k, MRR.
        \item \textbf{Evaluaci\'on online}: muestreo de respuestas, evaluaci\'on humana de fidelidad y relevancia.
        \item \textbf{M\'etricas de negocio}: tiempo ahorrado por analista, precisi\'on de decisiones.
    \end{itemize}
\end{frame}

\begin{frame}[fragile]{Prompting para minimizar alucinaciones}
    \begin{block}{Ejemplo de prompt}
        \begin{lstlisting}
Responde unicamente con la informacion de los documentos proporcionados.
Si la respuesta requiere combinar varios fragmentos, hazlo.
Cita los documentos al final de la respuesta.
Si no encuentras la informacion, di "No lo se".
        \end{lstlisting}
    \end{block}
    Incluir ejemplos few-shot mejora la adherencia.
\end{frame}

% ============================================================================
% SECCIÓN 10: CONEXIÓN CON EL RESTO DEL CURSO
% ============================================================================
\section{Conexi\'on con el resto del curso}

\begin{frame}{Relaci\'on con sesiones anteriores}
    \begin{itemize}
        \item \textbf{Sesi\'on 4 (Word Embeddings)}: base de los embeddings.
        \item \textbf{Sesi\'on 9 (Transformers)}: arquitectura de los LLMs usados en generaci\'on.
        \item \textbf{Sesi\'on 10 (BERT, GPT)}: modelos pre-entrenados que se usan en RAG.
        \item \textbf{Sesi\'on 11 (Tareas Avanzadas)}: QA y summarization se benefician de RAG.
    \end{itemize}
\end{frame}

% ============================================================================
% SECCIÓN 11: RESUMEN
% ============================================================================
\section{Resumen}

\begin{frame}{Lo que hemos aprendido}
    \begin{exampleblock}{Conceptos clave}
        \begin{itemize}
            \item \textbf{RAG}: arquitectura que combina recuperaci\'on y generaci\'on para superar limitaciones de LLMs.
            \item \textbf{Embeddings}: representaciones vectoriales para b\'usqueda sem\'antica.
            \item \textbf{Similitud del coseno}: m\'etrica para comparar vectores.
            \item \textbf{Bases de datos vectoriales}: FAISS, Pinecone, Weaviate, ChromaDB.
            \item \textbf{RAG Triad}: fidelidad, relevancia del contexto, relevancia de la respuesta.
            \item \textbf{RAG vs Fine-Tuning}: din\'amico vs est\'atico, menor costo y alucinaciones.
            \item \textbf{Implementaci\'on}: LangChain para prototipos r\'apidos.
            \item \textbf{Caso pr\'actico}: sistema de inteligencia de negocios.
        \end{itemize}
    \end{exampleblock}
\end{frame}

% --- Diapositiva final ---
\begin{frame}
    \centering
    \Huge \textbf{¡Gracias!}

\end{frame}

\end{document}